A set $\mathcal{F} \subset C(G, \Omega)$ is equicontinuous at a point $z_0$ in $G$ iff for every $\epsilon > 0$ there is a $\delta > 0$ such that for $|z - z_0| < \delta$,

$$d(f(z), f(z_0)) < \epsilon$$

for every $f$ in $\mathcal{F}$. $\mathcal{F}$ is equicontinuous over a set $E \subset G$ if for every $\epsilon > 0$ there is a $\delta > 0$ such that for $z$ and $z'$ in $E$ and $|z - z'| < \delta$,

$$d(f(z), f(z')) < \epsilon$$

for all $f$ in $\mathcal{F}$.

Notice that if $\mathcal{F}$ consists of a single function $f$ then the statement that $\mathcal{F}$ is equicontinuous at $z_0$ is only the statement that $f$ is continuous at $z_0$. The important thing about equicontinuity is that the same $\delta$ will work for all the functions in $\mathcal{F}$. Also, for $\mathcal{F} = \{f\}$ to be equicontinuous over $E$ is to require that $f$ is uniformly continuous on $E$. For a larger family $\mathcal{F}$ to be equicontinuous there must be uniform continuity.

Because of this analogy with continuity and uniform continuity the following proposition should not come as a surprise.